\chapter{Of the Deputy Mayor and the Order of Succession}
\label{art:succession}

\articlenote{In which is set down the office of the Deputy Mayor and the manner by which the silo chooses its next Mayor, that the office may never stand vacant and the silo's governance may never pass into uncertainty.}

\section{The Office of Deputy Mayor}

The Deputy Mayor is the second-highest office of the silo, subordinate to the Mayor in ordinary times. The Deputy Mayor's duty is to assist the Mayor in the conduct of the office and to hold in readiness the knowledge and authority to assume the Mayoralty should the sitting Mayor become unable to serve. The Deputy Mayor is appointed by the Mayor and confirmed by the Assembly under Article 20, serves at the Mayor's pleasure, and resides in the Up Top alongside the Mayor's chambers.

\begin{clauses}
\item The Deputy Mayor holds every duty of a citizen under this Pact and is subject to the same law as any other citizen, save where this Article or Article 4 expressly grants the Deputy some authority in the conduct of the Deputy's office.
\item The Deputy Mayor has access to all records, correspondence, and standing orders of the Mayor's office, that the Deputy may be prepared at any moment to assume the Mayoralty without loss of continuity.
\item No citizen afflicted with the Syndrome may hold the office of Deputy Mayor or be confirmed to it. Any Deputy Mayor who falls afflicted with the Syndrome shall immediately resign, and the Mayor shall appoint a new Deputy, confirmed by the Assembly under this Article.
\end{clauses}

\section{The Deputy Mayor's Assumption of the Mayoralty}

Upon the death of the sitting Mayor, or upon the Mayor's resignation, or upon the Mayor being removed from office by the Assembly under Article 20, the Deputy Mayor shall immediately assume the office and powers of the Mayor, and shall serve as Mayor until such time as a new Mayor is elected and confirmed under Section 3.

\begin{clauses}
\item The Deputy Mayor, upon assumption of the Mayoralty, becomes bound by every requirement and power set down in Article 4 for the Mayor, and Article 4 governs the Deputy's conduct from the moment of assumption forward.
\item The Deputy Mayor assuming the Mayoralty may issue standing orders, manage the silo's resources, and carry out all ordinary functions of the Mayor's office, save one: the Deputy may make no permanent appointment of a Department head or the Judge. Such permanent appointments are deferred until a new Mayor is elected and confirmed.
\item The Deputy Mayor's acts in the conduct of the Mayoralty during assumption are valid and binding on the silo and may not be reversed upon the election and assumption of a new Mayor, save where this Pact or a ruling of Judicial requires otherwise.
\item If the Mayor is absent but not dead or removed — taken ill, sent as an envoy, or otherwise prevented from the conduct of office for a known or suspected duration — the Deputy acts as Mayor with all powers including the power to make temporary appointments, such temporary authority ceasing upon the Mayor's return.
\end{clauses}

\section{The Election for the Next Mayor}

When the Mayoralty stands vacant by reason of death, resignation, or removal, an election shall be held. Any citizen of majority standing, with no debt to justice under Article 17 and no continuing sentence under Article 17, may stand for election as Mayor. The Assembly shall conduct the election and determine the process of the same.

\begin{clauses}
\item The Deputy Mayor, while serving as acting Mayor during the vacancy, may not stand for election to the Mayoralty, to ensure the impartiality of the acting office during the election itself.
\item A citizen elected to the Mayoralty must be confirmed by the Assembly under Article 20 before taking final office; confirmation shall not be withheld save for cause found by Judicial.
\item The newly elected and confirmed Mayor assumes office and all the powers of the Mayor under Article 4 upon confirmation; the Deputy Mayor's acting Mayoralty concludes at that moment.
\item The Assembly shall hold an election to choose and confirm the new Mayor; the Assembly, in consultation with the acting Mayor and the Heads of Departments, shall determine the timing and manner of the election.
\end{clauses}

\section{The Deputy Mayor's Election and Succession}

The Deputy Mayor is appointed by the sitting Mayor and confirmed by the Assembly under Article 20. Should the Deputy Mayor die, resign, or be removed from office before the Mayoralty stands vacant, the Mayor shall appoint a new Deputy, confirmed by the Assembly, to fill the vacancy.

\begin{clauses}
\item A new Deputy Mayor may be appointed at any time the sitting Mayor judges needful, and no more than one Deputy may serve at a time.
\item Should both the Mayor and the Deputy die or be removed simultaneously or in succession before a new Mayor can be elected, the highest-ranking Department head, determined by the order of their first confirmation in office, shall assume the Mayoralty as acting Mayor until the election is held and a new Mayor confirmed.
\end{clauses}
